\subsection{仿射空间的定义}

定义 1. 给定域 K 上的左（相应地，右）向量空间 T，相伴于 T 的仿射空间是加法群 T 的任意齐性空间 E（I，§5，第 5 号），使得 0 是 T 中唯一保持 E 的所有元素不变的算子（即 T 在 E 上忠实且传递地作用）。在这些条件下，T 称为 E 的平移空间，其元素称为 E 的平移（或 E 的自由向量）。

在下文中，我们将仅限于讨论 T 是 K 上的左向量空间的情形。仿射空间 E 的平移向量空间 T 的（在 K 上的）维数称为 E 的（在 K 上的）维数，记为 dim E 或 dim \(_{K}\) E。一维（相应地，二维）仿射空间称为仿射直线（相应地，仿射平面）。仿射空间的元素也称为点。

在定义 1 的条件下，对于 \(\mathbf{t} \in \mathbf{T}\) 与 \(a \in \mathbf{E}\) ，我们将用 \(\mathbf{t} + a\) 或 \(a + \mathbf{t}\) 表示点 \(a\) 在 \(\mathbf{t}\) 下的像。则关系式

\[
\mathbf {s} + (\mathbf {t} + a) = (\mathbf {s} + \mathbf {t}) + a, \quad 0 + a = a\tag{1}
\]

对 \(\mathbf{s} \in \mathrm{T}\) , \(\mathbf{t} \in \mathrm{T}\) , \(a \in \mathrm{E}\) 成立。映射 \(x \mapsto x + \mathbf{t}\) 是 \(\mathbf{E}\) 到其自身上的一个双射，我们把它与 \(\mathbf{t}\) 等同。此外，定义 1 还蕴含：对所有 \(a \in \mathrm{E}\) ，映射 \(\mathbf{t} \mapsto \mathbf{t} + a\) 是 \(\mathrm{T}\) 到 \(\mathbf{E}\) 上的一个双射。换言之，给定两个点 \(a\) , \(b\) （属于 \(\mathbf{E}\) ），存在且仅存在一个平移 \(\mathbf{t}\) 使得 \(b = \mathbf{t} + a\) ；我们把这个平移记作 \(b - a\) ；则公式

\[
a - a = 0, \quad a - b = - (b - a), \quad b = (b - a) + a\tag{2}
\]

\[
(c - b) + (b - a) = c - a
\]

对 \(a \in \mathbf{E}\) , \(b \in \mathbf{E}\) , \(c \in \mathbf{E}\) 成立。若四个点 \(a, b, a', b'\) （属于 \(\mathbf{E}\) ）满足 \(b - a = b' - a'\) ，则公式

\[
b ^ {\prime} = \left(b ^ {\prime} - b\right) + (b - a) + a = \left(b ^ {\prime} - a ^ {\prime}\right) + \left(a ^ {\prime} - a\right) + a
\]

以及 T 中加法的交换性表明 \(b' - b = a' - a\) .

给定一点 \(a \in E\) ，映射 \(x \mapsto x - a\) 是 \(E\) 到 \(T\) 上的一个双射；当 \(E\) 在此映射下与 \(T\) 等同时，我们就说把 \(E\) 视为通过取 \(a\) 作为 \(E\) 中的原点而得到的向量空间。反之，每个向量空间 \(T\) 都典范地具有一个相伴于 \(T\) 的仿射空间结构，即对应于子群 \(\{0\}\) （属于 \(T\) ）的齐性空间结构（I，§ 5，第 6 号）。

注记。本号中的定义以及后面的某些结果可立即推广到如下情形：我们不考虑向量空间 T，而是考虑任意一个带算子的交换群 T。

